\documentclass[10pt,a4paper]{article}
\usepackage[margin=1.8cm]{geometry}
\usepackage{multicol}
\usepackage{parskip}
\usepackage{titlesec}
\usepackage{enumitem}
\usepackage[hidelinks]{hyperref}
\usepackage{xcolor}
\usepackage{mdframed}

\definecolor{accent}{RGB}{160,90,0}
\definecolor{boxbg}{RGB}{255,248,230}
\titleformat{\section}{\large\bfseries\color{accent}}{}{0em}{}[\vspace{-4pt}\rule{\linewidth}{1pt}]
\titleformat{\subsection}{\normalsize\bfseries}{}{0em}{}
\setlist[itemize]{noitemsep, topsep=2pt, leftmargin=*}
\setlist[enumerate]{noitemsep, topsep=2pt, leftmargin=*}
\newmdenv[backgroundcolor=boxbg, linecolor=accent, linewidth=0.8pt, innerleftmargin=6pt, innerrightmargin=6pt, innertopmargin=4pt, innerbottommargin=4pt, skipabove=4pt, skipbelow=4pt]{keybox}

\begin{document}
\begin{center}
{\LARGE\bfseries\color{accent} Personal Development \& Productivity Reference}\\[4pt]
{\small Habits, Mindset, Focus, Learning, Mental Strength, Wellbeing, and Self-Discipline}
\end{center}
\vspace{4pt}\hrule\vspace{6pt}

\begin{multicols}{2}

\section{Habit Formation}

\subsection{Atomic Habits --- James Clear}
\textbf{4 Laws of Behaviour Change:}\\
Make it \textit{obvious} (cue) $\to$ make it \textit{attractive} (craving) $\to$ make it \textit{easy} (response) $\to$ make it \textit{satisfying} (reward).\\
\textbf{Key insight:} You do not rise to the level of your goals; you fall to the level of your systems.\\
\textbf{The 1\% better principle:} $1.01^{365} \approx 37$x improvement per year.\\
\textbf{Identity-based habits:} Decide the type of person you want to be; prove it with small wins.\\
\textbf{Habit stacking:} ``After I do X, I will do Y.''

\subsection{20 Micro Habits With a Massive Return on Life}
\begin{enumerate}
\item Wake up before 6\,am to reclaim quiet mornings.
\item Consume educational content at $1.5\times$ speed.
\item Give someone a genuine compliment each day.
\item Work on the most important project first thing.
\item Be heard more; offer fewer hot takes.
\item Fill a full glass of water as soon as you wake up.
\item Do 1 minute of stretching for every 30 minutes of sitting.
\item Write down your top 3 priorities for the day before starting.
\item Eat ``junk-free energy windows''; no challenging tasks in this time.
\item Learn to use the $4\times7\times8$ breathing technique to fall asleep quickly.
\item Save \emph{one} ``Subscription-like'' book on self-development to refer back to.
\item Course correct before you spiral; break the chain early.
\item Reach out to people just because they crossed your mind.
\item If you want to change something in your life, track it; you can't manage what you don't measure.
\item Put your phone in grayscale mode to reduce screen attractiveness.
\end{enumerate}

\subsection{10 Habits That Put You in the Top 1\%}
1. Protect your energy (guard attention).
2. Remove emotional reactions --- create a clear mind.
3. Start your day right (don't check phone immediately).
4. Be interested, not interesting (listen deeply).
5. Read 1 hour per day (30 books per year).
6. Simplify: break challenges into smallest possible steps.
7. Follow through completely.
8. Surround yourself with right people.
9. Track your progress.
10. Think long-term --- patience over instant gratification.

\subsection{10 Habits of People with an Elite Work Ethic}
1. Show up when it's inconvenient. 2. Set the bar higher. 3. They don't negotiate with laziness. 4. Track time ruthlessly. 5. They don't say ``I'll try.'' 6. They protect their momentum. 7. They move fast and fix later. 8. Work when others don't. 9. They embrace discomfort. 10. They don't wait for motivation.

\subsection{Daily Habits That Save You Time}
\begin{itemize}
\item Simplify your routine; fewer decisions = less fatigue.
\item Keep Essentials in Fixed Places (reduce search time).
\item Set Time Blocks for Work (batch similar tasks).
\item Reduce Phone Distractions (notifications off).
\item Follow the ``2-Minute Rule'': if it takes $\leq2$ min, do it now.
\end{itemize}

\subsection{5 Habits That Will Boost Your Brain Health}
1. \textbf{Japan's 80\% Rule:} Stop eating when 80\% full; eat slowly, avoid sugary drinks, hydrate well.
2. \textbf{Take walking meetings:} Steve Jobs, Nietzsche, Aristotle all did it.
3. \textbf{Have a strong rest ethic:} sleep is vital for memory consolidation.
4. \textbf{Avoid Energy Vampires:} spend time with people who energise you.
5. \textbf{Morning sunlight exposure:} 10--20 minutes resets circadian rhythm.

\section{Mental Strength and Mindset}

\subsection{10 Secrets to Become Mentally Unbreakable}
1. Don't like being alone $\to$ learn to love solitude.
2. Stop trying to control everything.
3. Do it now; don't overthink.
4. Trust the process even when nothing is happening.
5. Maintain a morning routine.
6. Embrace pain as a teacher.
7. Focus on what you CAN do, not what you can't.
8. Let go of the need for approval.
9. Be committed, not just interested.
10. Master your story --- the narrative you tell yourself.

\subsection{How to Be Mentally Strong}
1. Stop expecting from everyone. 2. Accept that life is unfair. 3. Don't beg for love or attention. 4. Keep emotions under control. 5. Learn to stay calm in chaos. 6. Stop taking things personally. 7. Walk away from toxic people. 8. Focus on solutions, not problems. 9. Believe in yourself always.

\subsection{Resilient Mindset (11 Rules)}
1. Don't expect it to be easy. 2. Separate facts from stories. 3. Focus on what you can change. 4. Take full responsibility. 5. Let go of what you can't change. 6. Learn from the past, plan for the future, live in the present. 7. Engage in physical and mental activity. 8. Embrace change and discomfort. 9. Connect with others. 10. Set boundaries. 11. Rest without guilt.

\subsection{12 Unstoppable Mindset Shifts}
``I have to do this'' $\to$ ``I \emph{get} to do this.''\\
``They're judging me'' $\to$ ``Their judgment doesn't change my value.''\\
``This could fail'' $\to$ ``This could change my life forever.''\\
``I feel stuck'' $\to$ ``I'm figuring it out.''\\
``I can't do this'' $\to$ ``I can figure this out if I take it one step at a time.''\\
``It takes too long'' $\to$ ``Every day I'm progressing.''\\
``I'm so exhausted'' $\to$ ``It'll take action first, the energy will follow.''\\
``I look stupid'' $\to$ ``I look like someone learning.''\\
``I'm not confident'' $\to$ ``Confidence comes from doing hard things.''\\
``I don't want to do this'' $\to$ ``I can't control the past, but I can control the next 5 minutes.''\\
``I've gone too far'' $\to$ ``I'm still in the game.''\\
``This is unfair'' $\to$ ``What can I do about it?''

\subsection{The Art of Discipline}
\begin{itemize}
\item Show up when you don't feel like it.
\item Do the work when motivation fades.
\item Keep promises you make to yourself.
\item Turn goals into habits.
\item Build confidence through small wins.
\end{itemize}
\textit{Discipline is not restriction. It is freedom earned through consistency.}

\section{Productivity and Focus}

\subsection{How to Be Seriously Organised}
1. Write down everything; don't trust your memory. 2. Use colour codes for categories. 3. Keep a running to-do list --- move today's tasks to your calendar. 4. Automate repetitive tasks. 5. Do a quick cleanup at the end of every day. 6. Batch similar tasks to avoid context-switching. 7. Lean on a folder system for digital and physical spaces. 8. Set similar tasks to stay on track for each task. 9. Stop multitasking. 10. Go paperless.

\subsection{Time Management: Where to Start?}
1. Identify how you currently spend your time.
2. Set your priorities.
3. Choose the one planning tool.
4. Write down your days and schedule.
5. Block time in your calendar.
6. Plan your day the night before (10--15 min).
7. Review your week on Sunday; reflect and plan for the coming week.

\subsection{How to Organise Your Brain (Strategic Thinking Framework)}
Plan first; apply second-order thinking; synthesis; problem-solving; abstraction; force multiplication; divergence; mental models. Long-term mindset: ask ``what is the 2nd-order consequence?''

\subsection{Overthinking vs Mindfulness}
\begin{itemize}
\item \textbf{Overthinking:} Stuck on past/future; catastrophising; analysing without action.
\item \textbf{Mindfulness:} Present moment; observing thoughts without judgement; accepting things as they are.
\end{itemize}
Mindfulness practice: 5-minute daily meditation, body scan, focused breathing (4-7-8 technique).

\subsection{The Research Loop}
\textit{The best researchers rarely involve a single stroke of genius. They involve relentless curiosity.}\\
Predict, Test, Update, Repeat. Read broadly. The pages in your head is how fast you learn. Descending reading speed matters; read papers with pen. Paranoia note: the first person you most avoid fooling is yourself.

\subsection{8 Trainable Skills Behind Better Research}
1. Pay your own intuition. 2. Upgrade your intuitions. 3. Tinker the loop. 4. Render on paper. 5. Work with your people. 6. Find your energy. 7. Go to the long game. 8. Know your goals.

\section{Learning and Memory}

\subsection{Improve Your Memory (12 Tips)}
1. Teach it. 2. Space repetition. 3. Create mnemonics. 4. Make it ordinary (link to existing knowledge). 5. Write it down (pen activates motor memory). 6. Say it out loud. 7. Chunk information (group into 3--5 items). 8. Use memory palaces (method of loci). 9. Engage senses (multi-modal encoding). 10. Recall actively (test yourself). 11. Don't multitask while learning. 12. Sleep well (consolidation happens during sleep).

\subsection{Study Smarter, Not Harder}
\begin{itemize}
\item \textbf{Active Recall:} Test yourself frequently; harder than re-reading but far more effective.
\item \textbf{Spaced Repetition:} Review material at increasing intervals (Anki, Leitner system).
\item \textbf{Practice Testing:} Simulate exam conditions; retrieval practice.
\item \textbf{Deep Work:} Focused distraction-free blocks (Cal Newport); 90-min sessions.
\item \textbf{Smart Note-Taking:} Zettelkasten method: atomic, linked notes.
\item \textbf{Consistency:} Small daily input beats marathon cram sessions.
\end{itemize}
``The Learning Formula:'' Learn $\to$ Recall $\to$ Master.

\subsection{Train Your Brain to Think Smarter}
Pause, Mental Subtraction, Compression, Zoom Out, Pattern Recognition, Environment Hacking. Practice cognitive flexibility daily.

\subsection{How Popular Programming Skills Take to Learn}
HTML/CSS: 2--3 months. Python/JavaScript: 6--12 months. Java/C/Go: 6--18 months. Rust: 12--24 months. Haskell: 18--36 months. For all: consistent daily practice beats irregular marathons. Build projects early.

\subsection{Your Brain Is Trainable}
1. What you keep doing, you keep.
2. Wake up, don't multitask.
3. Silence is powerful.
4. Saving \emph{tonight} builds tomorrow.
5. Saving habits build elite minds.
6. Own your focus. Start something.
7. Clear your space, clear your head.
8. Finish with cold water.
9. Do \emph{one thing} right now.

\subsection{The 3-Hour Window (Brain Optimal Performance)}
Your brain is in its peak state for ~3 hours per day. This window typically starts 30--60 min after waking, post coffee/breakfast. Use it exclusively for deep work: complex problem solving, writing, creative and strategic tasks. Never waste it on email, meetings, or admin.

\section{Wellbeing and Mental Health}

\subsection{8 Wellbeing Tips to Silence Anxiety}
1. Accept uncertainty as the natural state.
2. Practise box breathing (4-4-4-4).
3. Limit news and doomscrolling (set a 15-min daily limit).
4. Move your body (even a 20-min walk reduces cortisol).
5. Challenge catastrophic thinking (``What's the actual probability?'').
6. Connect with people who calm you.
7. Get enough sleep (7--9 h; anxiety magnifies with sleep deprivation).
8. Gratitude journalling: 3 things each morning.

\subsection{Habits vs Mindfulness (Book Reference Matrix)}
Atomic Habits (James Clear), The Power of Now (Eckhart Tolle), Habit Tools (Blinkist), 4 Mindset Books compared: Habit loops vs present-moment awareness, daily actions vs mental acceptance, identity shift vs mindset shift.

\subsection{Your Brain Lies (7 Habits That Rewire Your Thinking)}
1. Change how you talk to yourself.
2. Shift your state first.
3. Challenge the lies that lead to action.
4. Get triggers that lead to action.
5. Keep the reps small.
6. Interrupt old patterns.
7. Feed your brain new evidence.

\subsection{Neuroscience of Change}
Neuroplasticity: brain can reorganise throughout life. Change requires: discomfort (growth zone), repetition (myelination), emotion (stronger encoding), sleep (consolidation). Habits rewire neural pathways over ~66 days on average (not 21).

\subsection{10 Thinking Traps to Avoid}
1. Confirmation bias (seeking only confirming info). 2. Sunk cost fallacy. 3. Group-think. 4. Availability heuristic. 5. Anchoring bias. 6. Status quo bias. 7. Survivorship bias. 8. Attribution error. 9. Dunning-Kruger effect. 10. Recency bias.

\subsection{Dunning-Kruger Effect}
Peak of ``Mount Stupid'' (high confidence, low competence) $\to$ Valley of Despair $\to$ Slope of Enlightenment $\to$ Plateau of Sustainability. True experts have highest competence but are more aware of what they don't know.

\subsection{Brain Wave States}
Delta ($0.5$--$4$\,Hz): deep sleep, healing. Theta ($4$--$8$\,Hz): light sleep, creativity, subconscious. Alpha ($8$--$12$\,Hz): relaxed, focused, meditative. Beta ($12$--$30$\,Hz): active thinking, problem solving. Gamma ($30$--$100$\,Hz): high cognition, peak performance.

\subsection{How to Grow Your Intelligence}
1. Seek novelty: push into new domains regularly. 2. Challenge yourself: tasks slightly above current ability (flow state). 3. Think creatively: brainstorm with both sides of your brain. 4. Do things the hard way: struggle builds cognitive muscle. 5. Network: expose yourself to diverse thinkers.

\section{Motivation and Philosophy}

\subsection{Naval Mantra: ``Nobody claps for discipline''}
Nobody cheers for consistency. But one day, everyone will notice. Show up every day to fix your life. Do it alone. Do it broke. Do it tired. Do it scared. Just do it.

\subsection{12 Timeless Life Lessons (Japanese Wisdom)}
1. The quieter you become, the more you notice. 2. Let go of what no longer serves you. 3. A calm mind often solves more problems than a busy one. 4. Discipline brings peace; chaos creates stress. 5. Patience makes life gentler. 6. Suffering grows when you resist reality. 7. The less you carry, the further you go. 8. Holding onto anger often hurts you more than others. 9. Do simple things with complete attention. 10. Silence is a strength, not a weakness. 11. Ego shouts; wisdom whispers. 12. When you slow down, clarity emerges.

\subsection{Risk and Regret}
\begin{keybox}
Good advice is never as simple as ``Live for today'' or ``Save for the future.'' The only good advice is ``Minimise future regret.''
\end{keybox}

\subsection{``You may not see it yet, but everything is quietly aligning in your favour.''}
Trust the process.

\subsection{Don't Attach Yourself To:}
A person, a project, a company, an organisation.\\
\textbf{Attach yourself to:} A goal. A calling. A mission. A purpose.

\subsection{Quotes Worth Keeping}
``Intelligence is the ability to avoid doing work, yet getting the work done.'' --- Linus Torvalds\\
``Do not try to do everything. Do one thing well.'' --- Steve Jobs\\
``The measure of intelligence is the ability to change.'' --- Albert Einstein\\
``Comfort and luxury were once thought to be the chief requirements of life, when all that we need to make us happy is something to be enthusiastic about.'' --- Charles Kingsley\\
``Always assume things will work out, then do the work to make it true.'' --- A mentor

\subsection{``Feel the Fear; Do It Anyway''}
Don't stare straight at it; the fear will grow and it will disappear. The nature of fear is that people don't look at it. Thinking is important in some degree. If you accept the probabilistic thinking, you know that you are definitely wrong. Think failure is helpful in some degree, unless you are so close to something, you can see the actual consequences.

\section{Surround Yourself Well}

\subsection{Surround Yourself with Disciplined People}
Being around focused people minimises distractions and wasted time. You mirror the habits, ambitions, and standards of your circle. ``You are the average of the 5 people you spend the most time with.'' Proximity to discipline is contagious.

\subsection{How We Grow}
The fall $\to$ the bottom $\to$ the rise $\to$ the unstoppable. Why is this so difficult? ``I'm determined to stop'' (still fighting). Maybe I deserve this. ``I haven't discovered a different decision.'' ``I'm spiralling out of control.''

\end{multicols}
\end{document}
